\documentclass[11pt]{article}
\usepackage[a4paper,margin=27mm]{geometry}
\usepackage{amsmath,amssymb,amsthm,hyperref}
\newcommand{\per}{\operatorname{per}}
\newcommand{\OO}{\mathcal O}
\newcommand{\EE}{\mathcal E}
\newtheorem{theorem}{Theorem}
\title{A padding proof of a conjectured permanent inequality}
\author{}
\date{8 October 2026}
\begin{document}
\maketitle
\begin{abstract}
The unrestricted-cut permanent inequality stated as Conjecture~9.3 by Pan, Skandera, and Wang follows from their balanced inequality, Theorem~8.18, by adjoining an identity matrix. We give the complete reduction, including singular matrices and endpoint cuts. The only non-elementary input is the balanced theorem from their preprint.
\end{abstract}
For an integer $n\geq 0$, write $[n]=\{1,\dots,n\}$, let $\OO_n$ and $\EE_n$ be its odd and even elements, and write $A[S]=A_{S,S}$ for a principal submatrix in increasing index order. The empty permanent is $1$. A real matrix is \emph{totally nonnegative} (TNN) if every square minor is nonnegative.

We use the following even-order instance of Pan--Skandera--Wang's Theorem~8.18~\cite{PSW}: for every $m\geq2$ and every TNN matrix $B$ of order $2m$,
\begin{equation}\label{balanced}
\per B[[m]]\,\per B[\{m+1,\dots,2m\}]
\geq \per B[\OO_{2m}]\,\per B[\EE_{2m}].
\end{equation}
For $m=1$ this is an equality.

\begin{theorem}\label{main}
For every integer $n\geq1$, every TNN matrix $A\in\mathbb R^{n\times n}$, and every integer $0\leq h\leq n$,
\begin{equation}\label{goal}
\per A[[h]]\,\per A[\{h+1,\dots,n\}]
\geq\per A[\OO_n]\,\per A[\EE_n].
\end{equation}
In particular, Conjecture~9.3 of~\cite{PSW} holds.
\end{theorem}
\begin{proof}
First recall two elementary facts about ordered direct sums. If $C,D$ are TNN, then $C\oplus D$ is TNN: a square minor selecting unequal numbers of first-block rows and first-block columns is zero by rank deficiency, while a minor selecting equal numbers factors as the product of a minor of $C$ and a minor of $D$, without a sign change. Also,
\begin{equation}\label{factor}
\per(C\oplus D)=\per(C)\per(D),
\end{equation}
because only permutations preserving both blocks can contribute nonzero summands. These assertions include empty blocks. Every principal submatrix of an identity matrix has permanent $1$.

Suppose $0<h<n$ and $2h\leq n$. Put $d=n-2h$, $m=n-h$, and
\[
B=I_d\oplus A.
\]
This is a TNN matrix of order $2m$, with balanced leading and trailing principal blocks $I_d\oplus A[[h]]$ and $A[\{h+1,\dots,n\}]$. Thus the left side of~\eqref{balanced} equals the left side of~\eqref{goal}. An index $i$ of $A$ becomes $d+i$ in $B$. Consequently the two parity principal submatrices of $B$ are identity blocks direct-summed with $A[\OO_n]$ and $A[\EE_n]$, in the same order if $d$ is even and the opposite order if $d$ is odd. By~\eqref{factor},
\[
\per B[\OO_{2m}]\,\per B[\EE_{2m}]
=\per A[\OO_n]\,\per A[\EE_n].
\]
Applying~\eqref{balanced} proves the claim. The possible case $m=1$ is $n=2,h=1$, where the claim is an equality.

If instead $0<h<n$ and $2h\geq n$, put $d=2h-n$, $m=h$, and $B=A\oplus I_d$. The two balanced blocks are now $A[[h]]$ and $A[\{h+1,\dots,n\}]\oplus I_d$. The original indices of $A$ do not shift, and the parity principal submatrices simply acquire identity direct-summands. Thus both product equalities again hold, and~\eqref{balanced} applies. This also includes $d=0$.

Finally, if $h=0$ or $h=n$, the left side is $\per A$. Every entry of $A$ is nonnegative. The permutations preserving $\OO_n$ and $\EE_n$ contribute exactly the product on the right side, which is therefore at most the sum over all permutations. This handles the endpoint cuts and $n=1$.
\end{proof}

\paragraph{Scope and attribution.}
The proof uses no computation and no strict-positivity, irreducibility, or nonvanishing assumption. It is a short consequence of the balanced theorem in~\cite{PSW}, not an independent proof of that theorem. The cited version explicitly states the balanced result as Theorem~8.18 and the unrestricted-cut assertion as Conjecture~9.3. Since the balanced case is included in the unrestricted assertion, identity padding shows that the two universal statements are equivalent.

\begin{thebibliography}{1}
\bibitem{PSW}
S.~Pan, M.~Skandera, and J.~Wang,
\emph{Permanental inequalities and unit interval orders},
arXiv:2610.04809v1 [math.CO], 3 October 2026,
Theorem~8.18 and Conjecture~9.3.
\href{https://arxiv.org/abs/2610.04809v1}{arXiv:2610.04809v1}.
\end{thebibliography}
\end{document}
